\chapter{Loops \& Iteration}

In computer science, \textbf{Iteration} is the repetition of a process in order to generate a sequence of outcomes. Python provides two primary loop constructs: the `while` loop (condition-driven) and the `for` loop (sequence-driven).

\section{The `while` Loop}

A `while` loop repeatedly executes a block of code as long as a specified boolean condition remains `True`.

\begin{definitionbox}{Syntax of `while` Loop}
\begin{lstlisting}
count = 1
while count <= 5:
    print(f"Iteration: {count}")
    count += 1  # Increment to avoid an infinite loop!
\end{lstlisting}
\end{definitionbox}

\begin{videocard}{IITM BS Lecture: Introduction to While Loop}{https://www.youtube.com/watch?v=KTvVNN7ia8o}
Official lecture defining `while` loops, termination criteria, and preventing infinite loops.
\end{videocard}

\section{The `for` Loop and `range()` Function}

Python's `for` loop is designed to iterate directly over the elements of any sequence (such as a string, range, list, or tuple).

\begin{formulabox}{The `range(start, stop, step)` Function}
\begin{itemize}
    \item `range(n)` $\to$ generates integers from $0$ up to $n-1$.
    \item `range(a, b)` $\to$ generates integers from $a$ up to $b-1$.
    \item `range(a, b, s)` $\to$ generates integers from $a$ to $b-1$ stepping by $s$.
\end{itemize}
\end{formulabox}

\begin{examplebox}{Summing Numbers with a `for` Loop}
\begin{lstlisting}
n = 100
total_sum = 0

for i in range(1, n + 1):
    total_sum += i

print(f"Sum of 1 to 100: {total_sum}")  # Output: 5050
\end{lstlisting}
\end{examplebox}

\textbf{Data Science Relevance:} Loops allow data scientists to process rows of datasets, compute iterative optimization steps (like Gradient Descent in Machine Learning), and run Monte Carlo simulations.

\begin{videocard}{IITM BS Lecture: Introduction to For Loop}{https://www.youtube.com/watch?v=lvXuQ_x7EsI}
Official lecture explaining sequence iteration using the `for` loop.
\end{videocard}

\begin{videocard}{IITM BS Lecture: Range Function and Iteration}{https://www.youtube.com/watch?v=ihL3Eac9tFc}
Official lecture demonstrating `range()` start, stop, step parameters and iterating over strings directly.
\end{videocard}

\section{Nested Loops}

A loop can be placed inside another loop. The inner loop completes all its iterations for every single iteration of the outer loop.

\begin{examplebox}{Generating a Multiplication Table}
\begin{lstlisting}
for i in range(1, 4):
    for j in range(1, 4):
        print(f"{i}x{j}={i*j}", end="\t")
    print()  # New line after each row
\end{lstlisting}
\end{examplebox}

\begin{videocard}{IITM BS Lecture: Nested For Loops}{https://www.youtube.com/watch?v=-4MRaWABCuo}
Official lecture demonstrating multi-dimensional iteration using nested loops.
\end{videocard}

\section{Loop Control Statements: `break`, `continue`, and `pass`}

\begin{theorembox}{Loop Controls}
\begin{itemize}
    \item \textbf{`break`:} Immediately terminates the loop and jumps to the code outside.
    \item \textbf{`continue`:} Skips the rest of the current iteration and moves directly to the next iteration.
    \item \textbf{`pass`:} A null statement used as a placeholder when syntax requires a statement but no action is needed.
\end{itemize}
\end{theorembox}

\begin{videocard}{IITM BS Lecture: Break, Continue, and Pass}{https://www.youtube.com/watch?v=SVAVQHfJbE0}
Official lecture explaining loop control flow modification using `break`, `continue`, and `pass`.
\end{videocard}

\newpage
\section{Practice Exercises}
\textit{Detailed step-by-step solutions are available in the Appendix.}

\begin{enumerate}
    \item \textbf{Factorial Calculation:} Write a `while` loop that calculates the factorial of a number $N = 6$ ($6! = 6 \times 5 \times 4 \times 3 \times 2 \times 1 = 720$).

    \item \textbf{Pattern Printing:} Write nested `for` loops to print the following triangle pattern for $N = 4$ rows:
\begin{lstlisting}
*
**
***
****
\end{lstlisting}

    \item \textbf{Sum of Evens:} Write a `for` loop using `range()` that computes the sum of all \emph{even} numbers between $1$ and $50$ (inclusive).

    \item \textbf{Early Stopping (`break`):} Given a list of search scores `scores = [12, 45, 67, 89, 92, 40]`, write a `for` loop that iterates through scores and prints `"Threshold Met: X"` for the first score $\ge 80$, then immediately stops using `break`.

    \item \textbf{Data Science Simulation (Convergence Check):} In optimization, an algorithm updates a parameter $x$ until the difference between consecutive values is smaller than a tolerance $\epsilon = 0.001$. Simulate this using a `while` loop starting with `x = 10.0`, repeatedly updating `x = x * 0.5` until `x < 0.01`, counting the total iterations taken.
\end{enumerate}